\section{Introduction}
\label{sec:intro}

The European Union reached a cumulative PV capacity of 406~GW at the end of 2025, after adding 65.1~GW in that year \citep{SolarPowerEurope2025}. Belgium is among the most densely penetrated markets: its capacity approached 13~GW in 2025, roughly 1,080~W per capita \citep{IEAPVPS2026}. A large share of this capacity sits on residential rooftops. For the power system, these rooftops are a large population of small, sparsely documented generators whose orientation, tilt and even rated capacity are known only approximately.

%% [EPSR-SCOPE] power-system framing for Electric Power Systems Research; see JOURNAL_NOTES.md
Their output must be predicted a day in advance. \ifepsr Distribution system operators use day-ahead forecasts to anticipate reverse power flow and congestion on low-voltage feeders. Aggregators and balance-responsible parties bid distributed capacity into day-ahead markets and carry the imbalance cost of forecast error.\else Aggregators bid distributed capacity into day-ahead markets.\fi\ Home energy-management systems schedule overnight battery charging against expected generation. Each use needs generation at sub-hourly resolution for the following day, from a forecast issued around noon. At this horizon, numerical weather prediction (NWP) is in practice the only source of information about tomorrow's sky \citep{Inman2013,Antonanzas2016}.

Methods for PV forecasting have moved from persistence and autoregressive models to tree ensembles, deep sequence models and transformers \citep{Antonanzas2016,Yang2018}. A meta-analysis of 320 studies \citep{Nguyen2026} found that, at the day-ahead horizon, NWP input is worth +11.6 percentage points of skill score, about 1.6 times the largest model-family effect (+7.0 points for ensemble-hybrid models). The same work found that reported skill rose 3 points from 2006 to 2019 but 10 points from 2021 to 2022, which points to publication and selection bias rather than a step change in capability. If accuracy is set mainly by the input, then each input has a ceiling that no learner can exceed. Locating that ceiling needs an experiment that supplies information the operational forecast cannot have, not another benchmark comparison.

Three further problems are seldom quantified together on one dataset (Section~\ref{sec:related}). First, the \emph{reference problem}: without ground sensors, NWP irradiance is often validated against ERA5 \citep{Hersbach2020}, which is produced by the same centre and model as the forecast under test. Satellite retrievals such as SARAH-3 \citep{Pfeifroth2024} are an independent alternative. Second, the \emph{propagation problem}: model output statistics (MOS) correct NWP irradiance on the premise that better irradiance gives better power \citep{Verzijlbergh2015}. That premise holds for a physical model chain \citep{Mayer2021,Mayer2022}, but it is rarely tested when the downstream model is a learner trained on the same biased input. Third, the \emph{transfer problem}: most reported accuracy comes from chronological splits in which every system is also in the training set. This overstates what a newly connected rooftop receives \citep{Depoortere2025}, and the causes of transfer failure, as opposed to the average, are seldom examined.

This study addresses four research questions for a fleet of residential rooftops:
\begin{itemize}
\item[RQ1.] How much of the day-ahead error comes from the limits of the NWP input, and how much can the power model still address?
\item[RQ2.] Does widening the NWP input, through lagged re-runs, multi-model averaging or ensemble dispersion, move that limit?
\item[RQ3.] How much does the validation reference distort measured irradiance error, and does a correction learned against an independent reference improve power forecasts?
\item[RQ4.] What accuracy does a rooftop absent from training receive, what drives the failures, and how much local history is needed before site calibration helps?
\end{itemize}

The contributions are diagnostic rather than architectural:
\begin{enumerate}
\item A fleet-oracle experiment that bounds the value of perfect regional weather knowledge (RQ1).
\item Evidence that averaging across meteorological centres helps while averaging re-runs of one centre does not, and that ensemble dispersion adds no probabilistic skill to a quantile model (RQ2).
\item A measurement of reanalysis-reference distortion, checked against an independent ground-instrumented study, and an end-to-end test showing that a working irradiance correction does not propagate to power, on seen and unseen rooftops (RQ3).
\item A deployment benchmark for unseen rooftops: the full outcome distribution, a hyperparameter-selection protocol for transfer, rated-capacity error as the main failure mode, and the local history needed before calibration (RQ4).
\end{enumerate}
